\documentclass[exercice,a4]{thedocument}%
\usepackage{texmaths-tikz}%
\usepackage{pstricks}%
\usepackage{pst-bezier}%
\usepackage{pstricks-add}%
\usepackage{pst-eucl}%
\usepackage{pst-func}%
\usepackage{pst-tree}%
\usepackage{pst-circ}%
\usepackage{tikz}%
\usepackage{pgfplots}%
\usepackage{tkz-tab}%
\usepackage{tikzpagenodes}%
\usepackage{tabularray}%
\usepackage{tabularx}%
\usepackage{multirow}%
\usepackage{multicol}%
\usepackage{scratch3}%
\usepackage{enumitem}%
\usepackage[np]{numprint}%
\usepackage{pas-tableur}%
\usepackage{listings}%
\usepackage{graphicx}%
\usepackage{adjustbox}%
\usepackage{needspace}%
\usepackage{xcolor}%
\usepackage{eurosym}%
\usepackage{diagbox}%
\usepackage{colortbl}%
\usepackage{multido}%
\usepackage{icomma}%
\usepackage[french]{babel}%
\startdatakeys
    \source{26GENMATMEAG3}
    \BO{2026}
    \cycle{4}
    \competencesDuSocle{CA,RA,CO}
    \objectifsApprentissage{C4A1aa,C4A2ad,C4B2cb,C4D2aa,C4D3ac}
    \motsCles{programme de calcul,tableur,expression littérale,développement,factorisation,équation}
\enddatakeys
\definecolor{gris}{gray}{0.5}%
\providecommand{\wideparen}[1]{\overset{\frown}{#1}}%
\begin{document}%
\providecommand{\degres}{^\circ}%
On considère les deux programmes de calcul suivants.

\begin{center}
	\fbox{\begin{minipage}[t]{0.4\textwidth}
		{\centering \textbf{Programme A}\par}
		\begin{itemize}[nosep]
			\item Choisir un nombre.
			\item Lui ajouter 3.
			\item Calculer le carré de la somme obtenue.
			\item Soustraire 9 au résultat.
		\end{itemize}
	\end{minipage}}
	\hfill
	\fbox{\begin{minipage}[t]{0.48\textwidth}
		{\centering \textbf{Programme B}\par}
		\begin{itemize}[nosep]
			\item Choisir un nombre.
			\item Lui ajouter 6.
			\item Multiplier le résultat obtenu par le nombre de départ.
		\end{itemize}
	\end{minipage}}
\end{center}

\begin{enumerate}
	\item On choisit 4 comme nombre de départ.
	\begin{enumerate}
		\item Montrer qu'on obtient 40 avec le programme A.
		\item Quel résultat obtient-on avec le programme B ?
	\end{enumerate}
	\item On utilise un tableur pour obtenir d'autres résultats.

	\begin{center}
	\begin{tikzpicture}
		\tableur*[3]{A/4cm,B/1.2cm,C/1.2cm,D/1.2cm,E/1.2cm,F/1.2cm,G/1.2cm}
		\celtxt[width=4cm]{A}{1}{Nombre de départ}
		\celtxt[width=4cm]{A}{2}{Programme A}
		\celtxt[width=4cm]{A}{3}{Programme B}
		\celtxt*[width=1.2cm,align=center]{B}{1}{$-3$}
		\celtxt*[width=1.2cm,align=center]{C}{1}{$-1$}
		\celtxt[width=1.2cm,align=center]{D}{1}{1}
		\celtxt[width=1.2cm,align=center]{E}{1}{3}
		\celtxt[width=1.2cm,align=center]{F}{1}{7}
		\celtxt[width=1.2cm,align=center]{G}{1}{9}
		\celtxt*[width=1.2cm,align=center]{B}{2}{$-9$}
		\celtxt*[width=1.2cm,align=center]{C}{2}{$-5$}
		\celtxt[width=1.2cm,align=center]{D}{2}{7}
		\celtxt[width=1.2cm,align=center]{E}{2}{27}
		\celtxt[width=1.2cm,align=center]{F}{2}{91}
		\celtxt[width=1.2cm,align=center]{G}{2}{135}
		\selecCell{B}{3}
	\end{tikzpicture}
	\end{center}

	Quelle formule doit-on saisir dans la cellule B3 afin qu'elle soit étirée vers la droite pour compléter le tableau ?
\end{enumerate}

\begin{minipage}[t]{0.52\linewidth}
\begin{enumerate}[start=3]
	\item Soit $f$ la fonction qui, à un nombre de départ $x$, associe le nombre obtenu par le programme B. La courbe ci-contre représente la fonction $f$.
	\begin{enumerate}
		\item Déterminer graphiquement l'image de $-3$ par la fonction $f$.
		\item Déterminer graphiquement les antécédents de $-5$ par la fonction $f$.
	\end{enumerate}
	\item Déterminer par le calcul tous les nombres permettant d'obtenir 0 par le programme B.
	\item On désigne par $x$ le nombre choisi au départ.
	\begin{enumerate}
		\item Démontrer que : $(x + 3)(x + 3) = x^2 + 6x + 9$.
		\item Montrer que les deux programmes donnent le même résultat quel que soit le nombre $x$ choisi au départ.
	\end{enumerate}
\end{enumerate}
\end{minipage}\hfill
\begin{minipage}[t]{0.44\linewidth}
\vspace{0pt}
%Début Tikz
\begin{center}
	\begin{tikzpicture}%compiler avec pgfplots
		\begin{axis}[axis lines=middle,
			x={5mm}, y={5mm},
			xmin = -6.6, xmax = 3.2,
			ymin = -9.6, ymax = 6.3,
			xtick = {-6,...,3}, ytick = {-9,...,6},
			xticklabels = {}, yticklabels = {},
			grid=major,
			grid style={line width=1pt,gray!50},
			xlabel={$x$}, ylabel={$f(x)$},
			xlabel style={anchor=north east},
			ylabel style={anchor=north east},
			>=stealth, axis line style={line width=1pt,->}]
			\addplot [line width=1.2pt,color=red,smooth,samples=120,domain=-6.9:0.9]{x^2+6*x};
			\node[below right,font=\small] at (axis cs:0.8,0) {1};
			\node[left,font=\small] at (axis cs:0,1) {1};
			\node[below left,font=\small] at (axis cs:0,0) {0};
		\end{axis}
	\end{tikzpicture}
\end{center}
%Fin Tikz
\end{minipage}
\medskip
\showthesource
\end{document}
